\specsection{参考文献}

\begin{enumerate}
\def\labelenumi{(\Roman{enumi})}
\item
  Ch. DE LA VALLÉE POUSSIN, Intégrales de Lebesgue, Fonctions d'ensembles, Classes de Baire, Paris (Gauthier--Villars), 1st. edn., 1916; 2nd. edn., 1936.
\item
  M. FRÉCHET: a) Sur l'intégrale d'une fonctionelle étendue à un ensemble abstrait, Bull. Soc. Math. France, 43 (1915), pp.~248--265; b) Les familles et fonctions additives d'ensembles abstraits, Fund. Math., 4 (1923), pp.~329--365; ibid. 5 (1924), pp.~206-251.
\item
  S. SAKS, Theory of the Integral, 2nd. edn., New York (Stechert), 1937.
\item
  E. BOREL, Les probabilités dénombrables et leurs applications arithmétiques, Rend. Circ. Math. Palermo, 27 (1909), pp.~247-271.
\item
  H. STEINHAUS, Les probabilités dénombrables et leur rapport à la théorie de la mesure, Fund. Math., 4 (1923), pp.~286--310.
\item
  P. J. DANIELL: a) Integrals in an infinite number of dimensions, Ann. of Math., 20 (1918--19), pp.~281--288; b) Functions of limited variation in an infinite number of dimensions, ibid. 21 (1919--20), pp.~30--38.
\item
  B. JESSEN, The theory of integration in a space of an infinite number of dimensions, Acta Math., 63 (1934), pp.~249-323.
\item
  A. EINSTEIN, Investigations on the Theory of the Brownian Movement, New York (Dover), 1956.
\item
  P. LÉVY: a) Leçons d'Analyse Fonctionnelle, Paris (Gauthier-Villars), 1922（第二版于 1951 年由同一出版社以 Problèmes concrets d'Analyse Fonctionnelle 为题出版）; b) Processus stochastiques et mouvement brownien, Paris (Gauthier-Villars), 1948; c) Le mouvement brownien, Mémor. Sci. Math., 126 (1954).
\item
  N. WIENER, Differential space, J. Math. Phys. MIT, 2 (1923), pp.~131--174 (= Selecta, pp.~55--98, Cambridge (MIT Press), 1964).
\item
  R. E. A. C. PALEY 与 N. WIENER, Fourier transforms in the complex domain, Amer. Math. Soc. Colloq. Publ. No.~19, New York, 1934.
\item
  A. N. KOLMOGOROFF, Grundbegriffe der Wahrscheinlichkeit-srechnung, Berlin (Springer), 1933.
\item
  Ju. V. PROKHOROV, Convergence of random processes and limit theorems in probability theory, Theor. Probability Appl., 1 (1956), pp.~156--214.
\item
  S. BOCHNER, Harmonic Analysis and the Theory of Probability, Berkeley (University of California Press), 1960.
\item
  I. M. GELFAND: a) Generalized stochastic processes（俄文）, Dokl. Akad. Nauk SSSR, 100 (1955), pp.~853--856; b) Some problems of functional analysis（俄文）, Uspehi Mat. Nauk, 11 (1956), pp.~3--12 (= Amer. Math. Soc. Transl. (2), 16 (1960), pp.~315--324).
\item
  R. A. MINLOS, Generalized random processes and their extension to a measure（俄文）, Trudy Moskov. Math. Obšč., 8 (1959), pp.~497--518 (= Selected translations in mathematical statistics and probability, III (19), pp.~291--313).
\item
  I. M. GELFAND 与 N. Ya. VILENKIN, Generalized Functions, Vol. IV, New York (Academic Press), 1964（英译本）.
\item
  E. SPARRE-ANDERSSEN 与 B. JESSEN, On the introduction of measures in infinite product sets, Dansk. Vid. Selbskab. Mat. Fys. Medd., 25 (1948), No.~4, pp.~1-7.
\item
  A. D. ALEXANDROFF, Additive set functions in abstract spaces. I (Ch. 1), Mat. Sbornik, 8 (1940), pp.~307--348; II (Chs. 2 and 3), ibid., 9 (1941), pp.~563--628; III (Chs. 4 to 6), ibid., 13 (1943), pp.~169--238.
\item
  L. LE CAM, Convergence in distribution of stochastic processes, Univ. Cal. Publ. Statistics, No.~11 (1957), pp.~207--236.
\end{enumerate}

